% ── §11 Results: Gray Adjacency ──
\section{Results: Gray Adjacency Among Native Contacts}
\label{sec:results-gray}

\subsection{Motivation}
The Karnaugh-map framework predicts that residue pairs connected by a single
bit-flip in Gray-code space (Hamming distance $d_G = 1$) should be enriched
among native contacts, because the encoding is designed so that adjacent
codewords correspond to physicochemically similar amino acids.

\subsection{Result}
Among 47{,}430 native-contact pairs across 150 PSICOV proteins, the observed
fraction of pairs at Gray distance~1 is 22.83\% (95\% CI: [22.27\%, 23.40\%]
by protein-level bootstrap with 2{,}000 resamples).  The composition-controlled
background rate---constructed by within-protein, within-separation-bin
resampling of non-contact consensus residue pairs---is 19.37\%.  The resulting
enrichment of $1.178\times$ is highly significant by two-sided exact binomial
test ($p = 1.1 \times 10^{-77}$).

\subsection{Encoding Specificity}
To determine whether this enrichment is specific to the He/Gray encoding or
would arise from any clustered assignment of codewords to residues, we
performed a permutation test: we recomputed the enrichment under 30 random
relabelings of the 20 amino-acid codes while keeping all pair observations,
contact labels, and background construction fixed.  Random relabelings yield
a mean enrichment of $0.961 \pm 0.095$ (maximum 1.158), placing the true
He/Gray value ($1.198\times$) at $z = 2.51$ above the null distribution.
This confirms that the specific Gray structure carries genuine signal beyond
generic physicochemical clustering, although the majority of the raw
enrichment reflects the broader tendency of chemically similar residues to
contact each other.

\subsection{Comparison with Underpowered Analysis}
A previous analysis on the SARS-CoV-2 Spike protein alone found an identical
point estimate ($1.18\times$) but failed to reach significance ($p = 0.345$)
because only 38 contact pairs were available.  The present analysis, with
three orders of magnitude more data, resolves this ambiguity decisively:
the effect is real but small.
